\section{Countably generated Mittag-Leffler modules}
\label{section-ML-countable}

\noindent
It turns out that countably generated Mittag-Leffler modules have a
particularly simple structure.

\begin{lemma}
\label{lemma-ML-countable-colimit}
Let $M$ be an $R$-module.  Write $M = \colim_{i \in I} M_i$ where $(M_i,
f_{ij})$ is a directed system of finitely presented $R$-modules.  If $M$ is
Mittag-Leffler and countably generated, then there is a directed countable
subset $I' \subset I$ such that $M \cong \colim_{i \in I'} M_i$.
\end{lemma}

\begin{proof}
Let $x_1, x_2, \ldots$ be a countable set of generators for $M$.  For each $x_n$
choose $i \in I$ such that $x_n$ is in the image of the canonical map $f_i:
M_i \to M$; let $I'_{0} \subset I$ be the set of all these $i$.  Now
since $M$ is Mittag-Leffler, for each $i \in I'_{0}$ we can choose $j \in I$
such that $j \geq i$ and $f_{ij}: M_i \to M_j$ factors through
$f_{ik}: M_i \to M_k$ for all $k \geq i$  (condition (3) of Proposition
\ref{proposition-ML-characterization}); let $I'_1$ be the union of $I'_0$ with
all of these $j$.  Since $I'_1$ is a countable, we can enlarge it to a
countable directed set $I'_{2} \subset I$.  Now we can apply the same procedure
to $I'_{2}$ as we did to $I'_{0}$ to get a new countable set $I'_{3} \subset
I$.  Then we enlarge $I'_{3}$ to a countable directed set $I'_{4}$.  Continuing
in this way---adding in a $j$ as in Proposition
\ref{proposition-ML-characterization} (3) for each $ i \in I'_{\ell}$ if $\ell$
is odd and enlarging $I'_{\ell}$ to a directed set if $\ell$ is even---we get a
sequence of subsets $I'_{\ell} \subset I$ for $\ell \geq 0$.  The union $I' =
\bigcup I'_{\ell}$ satisfies:
\begin{enumerate}
\item $I'$ is countable and directed;
\item each $x_n$ is in the image of $f_i: M_i \to M$ for some $i
\in I'$;
\item if $i \in I'$, then there is $j \in I'$ such that $j \geq i$ and $f_{ij}:
M_i \to M_j$ factors through $f_{ik}: M_i \to M_k$ for all
$k \in I$ with $k \geq i$.  In particular $\Ker(f_{ik}) \subset \Ker(f_{ij})$
for $k \geq i$.
\end{enumerate}
We claim that the canonical map $\colim_{i \in I'} M_i \to
\colim_{i
\in I} M_i = M$ is an isomorphism.  By (2) it is surjective.  For injectivity,
suppose $x \in \colim_{i \in I'} M_i$ maps to $0$ in $\colim_{i \in
I} M_i$.
Representing $x$ by an element $\tilde{x} \in M_i$ for some $i \in I'$, this
means that $f_{ik}(\tilde{x}) = 0$ for some $k \in I, k \geq i$.  But then by
(3) there is $j \in I', j \geq i,$ such that $f_{ij}(\tilde{x}) = 0$.  Hence $x
= 0$ in $\colim_{i \in I'} M_i$.
\end{proof}

\noindent
Lemma \ref{lemma-ML-countable-colimit}
implies that a countably generated Mittag-Leffler module $M$ over
$R$ is the colimit of a system
$$
M_1 \to M_2 \to M_3 \to M_4 \to \ldots
$$
with each $M_n$ a finitely presented $R$-module. To see this argue as in the
proof of
Lemma \ref{lemma-ML-limit-nonempty}
to see that a countable directed set has a cofinal
subset isomorphic to $(\mathbf{N}, \geq)$. Suppose
$R = k[x_1, x_2, x_3, \ldots]$ and $M = R/(x_i)$. Then $M$ is finitely
generated but not finitely presented, hence not Mittag-Leffler (see
Example \ref{example-ML} part (1)).
But of course you can write $M = \colim_n M_n$ by taking
$M_n = R/(x_1, \ldots, x_n)$, hence the condition that you can write
$M$ as such a limit does not imply that $M$ is Mittag-Leffler.

\begin{lemma}
\label{lemma-ML-countable}
Let $R$ be a ring.
Let $M$ be an $R$-module.
Assume $M$ is Mittag-Leffler and countably generated.
For any $R$-module map $f : P \to M$ with $P$ finitely generated there
exists an endomorphism $\alpha : M \to M$ such that
\begin{enumerate}
\item $\alpha : M \to M$ factors through a finitely presented $R$-module, and
\item $\alpha \circ f = f$.
\end{enumerate}
\end{lemma}

\begin{proof}
Write $M = \colim_{i \in I} M_i$ as a directed colimit of finitely
presented $R$-modules with $I$ countable, see
Lemma \ref{lemma-ML-countable-colimit}.
The transition maps are denoted $f_{ij}$ and we use $f_i : M_i \to M$
to denote the canonical maps into $M$. Set $N = \prod_{s \in I} M_s$. Denote
$$
M_i^* = \Hom_R(M_i, N) = \prod\nolimits_{s \in I} \Hom_R(M_i, M_s)
$$
so that $(M_i^*)$ is an inverse system of $R$-modules over $I$.
Note that $\Hom_R(M, N) = \lim M_i^*$.
As $M$ is Mittag-Leffler, we find for every
$i \in I$ an index $k(i) \geq i$ such that
$$
E_i := \bigcap\nolimits_{i' \geq i} \Im(M_{i'}^* \to M_i^*)
=
\Im(M_{k(i)}^* \to M_i^*)
$$
Choose and fix $j \in I$ such that $\Im(P \to M) \subset \Im(M_j \to M)$.
This is possible as $P$ is finitely generated. Set $k = k(j)$.
Let
$x = (0, \ldots, 0, \text{id}_{M_k}, 0, \ldots, 0) \in M_k^*$ and
note that this maps to $y = (0, \ldots, 0, f_{jk}, 0, \ldots, 0) \in M_j^*$.
By our choice of $k$ we see that $y \in E_j$. By
Example \ref{example-ML-surjective-maps}
the transition maps $E_i \to E_j$ are surjective for each $i \geq j$
and $\lim E_i = \lim M_i^* = \Hom_R(M, N)$. Hence
Lemma \ref{lemma-ML-limit-nonempty}
guarantees there exists an element $z \in \Hom_R(M, N)$
which maps to $y$ in $E_j \subset M_j^*$. Let $z_k$ be the $k$th component
of $z$. Then $z_k : M \to M_k$ is a homomorphism such that
$$
\xymatrix{
M \ar[r]_{z_k} & M_k \\
M_j \ar[ru]_{f_{jk}} \ar[u]^{f_j}
}
$$
commutes. Let $\alpha : M \to M$ be the composition
$f_k \circ z_k : M \to M_k \to M$.
Then $\alpha$ factors through a finitely presented module by construction and
$\alpha \circ f_j = f_j$. Since the image of $f$ is contained in the image of
$f_j$ this also implies that $\alpha \circ f = f$.
\end{proof}

\noindent
We will see later (see
Lemma \ref{lemma-split-ML-henselian})
that
Lemma \ref{lemma-ML-countable}
means that a countably generated Mittag-Leffler module over a henselian local
ring is a direct sum of finitely presented modules.




